\documentclass[11pt]{article}

\usepackage[T1]{fontenc}
\usepackage{lmodern}
\usepackage{amsmath,amssymb,amsthm,mathtools}
\usepackage{microtype}
\usepackage[margin=1in]{geometry}
\usepackage[colorlinks=true,linkcolor=blue,citecolor=blue,urlcolor=blue]{hyperref}

\newtheorem{theorem}{Theorem}
\newtheorem{lemma}{Lemma}
\newtheorem{proposition}{Proposition}
\newtheorem{corollary}{Corollary}
\theoremstyle{definition}
\newtheorem{definition}{Definition}
\theoremstyle{remark}
\newtheorem{remark}{Remark}

\title{The Erd\H{o}s--Gy\'arf\'as Balanced-Coloring Conjecture\\
for Twelve Colors}
\author{Author information to be finalized}
\date{\textbf{Working draft — not for submission}\\31 October 2026}

\begin{document}
\maketitle

\begin{abstract}
We prove the fixed twelve-color case of the Erd\H{o}s--Gy\'arf\'as
balanced-coloring conjecture: there is no edge-coloring of $K_{145}$ with
twelve colors for which every thirteen-vertex set contains an edge of every
color. The proof selects a color class with the fewest edges and develops a
hereditary extremal-graph reduction. A finite recurrence reduces a
hypothetical counterexample to a sixty-one-vertex residual graph in a narrow
edge range. The terminal analysis combines low-edge obstructions on
forty-seven and forty-eight vertices, a repivot/charge argument, and a
uniform terminal theorem. The proof is accompanied by finite computational
certificates and an independent AI-assisted informal audit of a pinned
source snapshot. The argument is not yet formalized in a proof assistant, and
external human peer review remains to be completed.
\end{abstract}

\section{Introduction}

A coloring of the edges of a complete graph with $r$ colors is called
\emph{balanced} if every set of $r+1$ vertices spans at least one edge of
each color. Erd\H{o}s and Gy\'arf\'as formulated the conjecture that
$K_{r^2+1}$ has no balanced $r$-edge-coloring for every integer $r\geq 3$
\cite{ErdosGyarfas1999}. Their paper proves the cases $r=3$ and $r=4$.
The five-color case was subsequently established by Xia and Szeider
\cite{XiaSzeider2026}, using a combination of formalized reduction,
SAT-based enumeration, and independently checkable certificates.

The purpose of this paper is to establish the next fixed-parameter statement
addressed by this project.

\begin{theorem}[The twelve-color case]
\label{thm:main}
There is no edge-coloring
\[
  c:E(K_{145})\longrightarrow \{1,2,\ldots,12\}
\]
such that every $13$-vertex subset of $K_{145}$ contains an edge of every
one of the twelve colors.
Equivalently, every $12$-edge-coloring of $K_{145}$ has a $13$-vertex set
whose induced edges omit at least one color.
\end{theorem}

Theorem~\ref{thm:main} is a fixed-$r$ result. It does not settle the
conjecture for all $r\geq 3$.

\subsection{Overview of the proof}

Assume for contradiction that a balanced twelve-coloring of $K_{145}$ exists.
Choose a color class with the fewest edges and regard it as a graph $G$.
The balanced-coloring hypothesis imposes hereditary upper bounds on the
number of edges induced by every vertex subset of $G$. These bounds are
combined with classical extremal graph results, including Tur\'an-type
bounds, Brooks' theorem, the Andr\'asfai--Erd\H{o}s--S\'os theorem, and a
second-Tur\'an bound for non-partite $K_{r+1}$-free graphs
\cite{Turan1941,Brooks1941,AndrasfaiErdosSos1974,KangPikhurko2005}.

The front-end reduction first rules out a contextual $R(4,49)$ residual
inside the original full-coloring setting. A packing and recurrence argument
then leaves a unique sixty-one-vertex residual $F$ with
\[
  397\leq e(F)\leq 408,\qquad
  \alpha(F)\leq 5,\qquad \omega(F)\leq 11,
\]
together with the required hereditary induced-subgraph caps.

The terminal argument separates according to the minimum degree of $F$.
The degree-twelve branch is excluded using a forty-eight-vertex
310-edge lower bound and an inserted-center contradiction. The degree-thirteen
branch is handled by degree-thirteen pivots, exact repivot identities, and a
uniform charge theorem. A forty-seven-vertex low-edge theorem supplies the
287-edge core floor, while the terminal charge inequalities rule out the
remaining edge values.

\subsection{Organization}

Section~\ref{sec:reduction} records the coloring reduction and the
sixty-one-vertex terminal contract. Section~\ref{sec:cores} states the
forty-seven- and forty-eight-vertex core theorems. Section~\ref{sec:terminal}
describes the h12 and h13 terminal arguments. Section~\ref{sec:computation}
specifies the finite computational verification and its trust boundary.
The detailed proof transcription and all exact hypotheses are being
completed from the frozen source documents identified in the repository's
publication index.

\section{Notation and the hereditary graph contract}

All graphs in this paper are finite and simple. For a graph $H$, write
$v(H)$ and $e(H)$ for its numbers of vertices and edges, $\delta(H)$ for
its minimum degree, $\alpha(H)$ for its independence number, and
$\omega(H)$ for its clique number. For $U\subseteq V(H)$, $H[U]$ denotes
the induced subgraph.

The final proof uses hereditary induced-subgraph edge caps inherited from
the original coloring assumption. Their exact definition and the derivation
from the least-edge color class must be included in the final proof
transcription; they are not additional assumptions on an arbitrary graph.

\section{The coloring reduction and the T61 residual}
\label{sec:reduction}

\paragraph{Source of the proof.}
The complete argument for this section is the frozen document
\texttt{K145\_TO\_T61\_FRONTEND\_AUDIT.md} at integrated proof commit
\texttt{c5f05f5868ed4d1347dc6d16f9d5b5be230ac7e2}. It proves the
color-class caps, contextual $R(4,49)$ exclusion, exact packing/edge ledger,
all recurrence rules, and the unique T61 stop.

\paragraph{Terminal contract.}
The recurrence has 1,752 nodes and 144 outer-stop comparisons. Its unique
surviving stop has parameters $d=11$ and $j=6$ and yields a graph $F$ on
61 vertices with $397\leq e(F)\leq408$, $\alpha(F)\leq5$,
$\omega(F)\leq11$, and all inherited induced-subgraph caps.

\paragraph{Proof transcription status.}
The complete line-by-line derivation of the cap function, packing
transversal, edge partition, recurrence seeds, and matching bounds is being
transcribed from the frozen proof source. This paragraph is a roadmap, not a
replacement for those derivations.

\section{The forty-seven- and forty-eight-vertex core theorems}
\label{sec:cores}

The integrated proof establishes a universal forty-seven-vertex low-edge
obstruction. Under the exact contract
\[
  v(H)=47,\quad \delta(H)\geq12,\quad
  \alpha(H)\leq4,\quad \omega(H)\leq11,
\]
there is no such graph with $e(H)\leq286$. In particular, every
forty-seven-vertex core satisfying the corresponding Ramsey--Tur\'an
contract has at least 287 edges.

A separate forty-eight-vertex argument proves the 310-edge floor required
for the degree-twelve terminal. This argument uses its own minimum-degree
gate and a complete charge-$20$ shell classification; it is not obtained
by transferring the forty-seven-vertex bound to a different order.

The exact hypotheses, equality structures, common-neighborhood capacity,
shared-excess accounting, and support-orbit classification must be included
in the final proof text from the frozen low-edge, equality, and h12 sources.

\section{The sixty-one-vertex terminal}
\label{sec:terminal}

The minimum-degree cases are separated as follows.

\begin{itemize}
  \item If $\delta(F)\leq 11$, the original full-color context yields a
  forbidden contextual $R(4,49)$ object.
  \item If $\delta(F)\geq14$, the degree sum forces more than 408 edges.
  \item If $\delta(F)=12$, the forty-eight-vertex floor and the complete
  charge-$20$ equality classification lead to the inserted-center
  contradiction.
  \item If $\delta(F)=13$, every regular pivot produces a forty-seven-point
  core. The core has at least 287 edges, and the COVER charge theorem
  excludes the remaining energies 407 and 408.
\end{itemize}

For a degree-thirteen pivot $p$, the proof uses an exact identity of the form
\[
  e(C_p)=E-91-q_p,
\]
where $C_p=F-N_F[p]$ and $q_p$ is the pivot charge. The full definitions,
repivot identities, quantifier order over all regular pivots, and terminal
classification are part of the frozen COVER and h13 proof sources and must
be transcribed before submission.

\section{Finite verification and reproducibility}
\label{sec:computation}

The independent audit reports the following finite checks for the pinned
integrated proof:

\begin{center}
\begin{tabular}{lr}
\hline
Obligations in the audit matrix & 58\\
Q27 / Q28 canonical states & 2 / 8\\
Subset identities replayed & 81,920\\
Recurrence nodes replayed & 1,752\\
Outer-stop comparisons & 144\\
\hline
\end{tabular}
\end{center}

The finite computations support bounded classifications and arithmetic
replay after the analytic reductions. They do not replace the analytic
proofs of the graph-to-state mappings or their completeness. The independent
replay is implemented in Python 3.12.14 using the standard library and is
documented in \texttt{PUBLICATION/REPRODUCIBILITY.md}.

\section{Scope, limitations, and formalization status}

The theorem proved here is the fixed $r=12$ instance. The proof does not
claim to settle the general conjecture. It also does not claim that the
abstract uncolored high-energy $R(4,49)$ layers are empty; the proof instead
uses a contextual contradiction available inside the original full-color
setting.

The current proof is an informal mathematical proof with finite
computational verification and an independent AI-assisted audit. It has
not yet been formalized in Lean 4/Mathlib and has not yet undergone external
human peer review. These qualifications will remain explicit in any
preprint version.

\bibliographystyle{plain}
\bibliography{references}

\end{document}
